\makeabstract
{ % #1: Title
Essential Roles of Six2 in Cranial Base Synchondrosis Development
}
{ % #2: Authorship
Aranne Jung\textsuperscript{1, 2},
Yuta Nakai\textsuperscript{2},
Azusa Shimizu\textsuperscript{2},
Noriaki Ono\textsuperscript{2}
}
{ % #3: Affiliations (use \\ to start a newline)
\textsuperscript{1} Amherst College, Amherst, MA\\
\textsuperscript{2} Department of Diagnostic and Biomedical Sciences, UTHealth Houston School of Dentistry, Houston, TX
}
{ % #4: Purpose
The cranial base is a skeletal structure that contributes to skull growth and elongation. The cranial base synchondroses are cartilaginous joints and sites of endochondral growth, with the spheno-occipital synchondrosis (SOS) and intersphenoid synchondrosis (ISS) serving as major growth sites. Unlike long-bone growth plates, skeletal progenitor populations in the cranial base synchondroses remain poorly characterized. \textit{Six2} is a transcription factor essential for the development of the kidney, limbs, and eyes. Single-cell RNA sequencing has identified \textit{Six2} as highly expressed within the cranial base synchondroses. Global \textit{Six2} knockout in mice results in premature SOS ossification and ISS loss, suggesting a crucial role for \textit{Six2} in cranial base synchondrosis development. This project aimed to investigate the functional importance of \textit{Six2} in postnatal cranial base synchondrosis development.
}
{ % #5: Methods
Using the \textit{Cre-loxP} system, \textit{Six2} was conditionally inactivated in postnatal chondrocytes using \textit{Col2a1-CreER; Six2\textsuperscript{fl/fl}; R26R\textsuperscript{tdTomato}} (Col2a1-Six2-cKO) mice with tamoxifen injection at postnatal day 3 (P3). We also conditionally inactivated the chondrocyte master regulator \textit{Sox9} in \textit{Six2\textsuperscript{+}} cells using \textit{Six2-creER; Sox9\textsuperscript{fl/fl}; R26R\textsuperscript{tdTomato}} (Six2-Sox9-cKO) mice by injecting tamoxifen at P3.
}
{ % #6: Results
No significant differences in SOS or ISS length were observed between Col2a1-Six2-cKO and their controls. Lineage-tracing with the tdTomato reporter showed comparable labeling levels between Col2a1-Six2-cKO and control samples. In contrast, Six2-Sox9-cKO mice exhibited a significantly truncated ISS length associated with apparent premature ossification; however, interestingly, no significant difference in SOS length was observed. Increased tdTomato\textsuperscript{+} cells in Six2-Sox9-cKO mice suggest that \textit{Sox9} loss in \textit{Six2\textsuperscript{+}} cells may affect chondrocyte differentiation.
}
{ % #7: Conclusion
Overall, these findings suggest that \textit{Sox9} is required for \textit{Six2\textsuperscript{+}} cells to differentiate into chondrocytes and maintain postnatal cranial base synchondrosis particularly in ISS, while \textit{Six2} itself is not required in these chondrocytes. Future studies will investigate the relationship between \textit{Six2} and other chondrogenic regulatory genes to understand its role in cranial base synchondrosis development.
}
 \newpage